\section*{الفصل \ref{ch:g12:oscillators-and-time} --- المذبذبات الميكانيكية وقياس الزمن}

\begin{solution}{exo:g12:oscillators-and-time:1}
(أ) $T_0 = 2\pi\sqrt{1.00/9.81} = \qty{2.01}{s}$؛
$f_0 = 1/T_0 = \qty{0.50}{Hz}$. (ب) $T_0 \propto \sqrt{L}$: فيتضاعف
\omterm{def:g12:oscillators-and-time:oscillator}{الدور}.
\end{solution}

\begin{solution}{exo:g12:oscillators-and-time:2}
القلب: $f = 72/60 = \qty{1.2}{Hz}$، $T = \qty{0.83}{s}$. والجناح:
$T = 1/30 = \qty{33}{ms}$. وارتفاع الصدر هو \emph{\omterm{def:g12:oscillators-and-time:oscillator}{السعة}}.
\end{solution}

\begin{solution}{exo:g12:oscillators-and-time:3}
$T_0 = 2\pi\sqrt{0.10/25} = \qty{0.40}{s}$، $f_0 = \qty{2.5}{Hz}$.
ومضاعفة $m$: $T_0 \times \sqrt{2} \approx 1.4$.
\end{solution}

\begin{solution}{exo:g12:oscillators-and-time:4}
\omterm{def:g12:oscillators-and-time:damping}{شبه الدور} \qty{2.0}{s}. والنسب $1.5/2.0 = 1.13/1.5 = 0.85/1.13
= 0.75$ في كل \omterm{def:g12:oscillators-and-time:oscillator}{دور}؛ والعظمى التالية $0.85 \times 0.75 \approx
\qty{0.64}{cm}$.
\end{solution}

\begin{solution}{exo:g12:oscillators-and-time:5}
النوّاس: $T_0 \propto 1/\sqrt{g}$، ومنه $\times\sqrt{9.81/1.62} =
2.46$ (أي أبطأ). \omterm{prop:g12:oscillators-and-time:spring-period}{ومذبذب الزنبرك والكتلة}: لا يتغير --- إذ لا $g$ في
$2\pi\sqrt{m/k}$.
\end{solution}

\begin{solution}{exo:g12:oscillators-and-time:6}
(أ) $[L/g] = \unit{m}/(\unit{m/s^2}) = \unit{s^2}$؛ وجذره التربيعي
زمن. (ب) ولا يظهر \omterm{def:g10:orders-of-magnitude:unit}{الكيلوغرام} إلا في $m$: فلا شيء آخر يقدر أن
يلغيه، ومن ثم لا تقدر $m$ أن تدخل. (ج) ولا أبدًا: فالأعداد الخالصة مثل $2\pi$
غير مرئية للأبعاد.
\end{solution}

\begin{solution}{exo:g12:oscillators-and-time:7}
$\frac{dx}{dt} = -\frac{2\pi}{T_0}A\sin(2\pi t/T_0 + \varphi)$، ومنه
$\frac{d^2x}{dt^2} = -(2\pi/T_0)^2 x(t)$. والمساواة مع $-(k/m)x$
من أجل كل $t$ تصحّ بالضبط حين $(2\pi/T_0)^2 = k/m$، أي\
$T_0 = 2\pi\sqrt{m/k}$.
\end{solution}

\begin{solution}{exo:g12:oscillators-and-time:8}
$E_m = \tfrac12 kA^2 = \tfrac12 \times 40 \times 0.030^2 =
\qty{1.8e-2}{J}$؛ $v_{\max} = \sqrt{2E_m/m} = \sqrt{0.18} =
\qty{0.42}{m/s}$ --- عند $x = 0$؛ وتنعدم عند الطرفين $x = \pm A$.
\end{solution}

\begin{solution}{exo:g12:oscillators-and-time:9}
(أ) $L = gT_0^2/4\pi^2 = \qty{0.994}{m}$. (ب) $\Delta T_0/T_0 =
\tfrac12 \times 1.0\% = 0.5\%$: $86400 \times 0.005 \approx
\qty{4.3e2}{s}$ --- أي سبع دقائق تأخّرًا في اليوم.
\end{solution}

\begin{solution}{exo:g12:oscillators-and-time:10}
(أ) $A \propto \sqrt{E}$: $\sqrt{0.95} \approx 0.975$، أي هبوط قدره $2.5\%$.
(ب) $0.95^n = 0.5$: $n = \ln 2/\ln(1/0.95) \approx 14$
دورًا. (ج) ولا: ففي \omterm{def:g12:oscillators-and-time:damping}{التخامد} الخفيف يبقى في حدود جزء من
المئة من $T_0$.
\end{solution}

\begin{solution}{exo:g12:oscillators-and-time:11}
(أ) عند $f_0$ تصل كل دفعة على وتيرة الحركة فتضيف \omterm{def:g10:energy-conservation:energy}{طاقة}؛ وبعيدًا عن الذروة
تساعد الدفعات وتعوق بالتناوب. (ب) والمشي على الخطى إثارة
\omterm{def:g10:signals-and-waves:period}{دورية}؛ وقرب $f_0$ الخاص بالجسر كانت ستتسلق ذروة
\omterm{def:g12:oscillators-and-time:resonance}{الرنين}. (ج) وتمسح الحلّة \omterm{def:g12:oscillators-and-time:resonance}{التواتر الذاتي} للهيكل
أثناء تسارع الدوران: \omterm{def:g12:oscillators-and-time:resonance}{فرنين} لحظي، ثم تجاوز للذروة.
\end{solution}

\begin{solution}{exo:g12:oscillators-and-time:12}
(أ) $T_0 \propto 1/\sqrt{g}$: أي بطيئة، بمقدار $\Delta T_0/T_0 = \tfrac12
\times 0.03/9.81 = \num{1.5e-3}$، أي\ $\approx \qty{132}{s}$ في
اليوم. (ب) و$L \propto g$ عند $T_0$ ثابتة: فقصّرها بمقدار $0.994 \times
0.03/9.81 \approx \qty{3.0}{mm}$.
\end{solution}

\begin{solution}{exo:g12:oscillators-and-time:13}
$v_{\max} = 2\pi A/T_0 = 2\pi \times 0.040/0.63 = \qty{0.40}{m/s}$؛
$a_{\max} = (2\pi/T_0)^2 A = \qty{4.0}{m/s^2}$؛ $k = m(2\pi/T_0)^2
= \qty{20}{N/m}$؛ $E_m = \tfrac12 kA^2 = \qty{1.6e-2}{J}$.
\end{solution}

\begin{solution}{exo:g12:oscillators-and-time:14}
$T_0 = 40.1/20 = \qty{2.005}{s}$؛ $g = 4\pi^2 L/T_0^2 = 4\pi^2
\times 0.995/4.020 = \qty{9.77}{m/s^2}$. والارتياب $\approx 0.2\%
+ 2 \times 0.5\% = 1.2\%$، أي\ $\pm\qty{0.12}{m/s^2}$: \omterm{def:g11:electric-gravitational-fields:field}{فالمجال}
يحوي \qty{9.81}{m/s^2} --- وهو متسق.
\end{solution}

\begin{solution}{exo:g12:oscillators-and-time:15}
النوّاس $10/86400 \approx \num{1.2e-4}$؛ والكوارتز $0.5/\num{2.6e6}
\approx \num{1.9e-7}$؛ والسيزيوم $1/\num{3.2e15} \approx \num{3e-16}$.
وخطأ قدره \qty{1}{\micro s} يعطي $\num{1000} \times \qty{30}{cm} =
\qty{300}{m}$ في الموضع: فنظام GPS يحتاج العائلة الذرّية.
\end{solution}

\begin{solution}{pb:g12:oscillators-and-time:1}
\textbf{1.} يُقسم خطأ ردّ الفعل $\approx \qty{0.2}{s}$ على
$50$: أي نحو \qty{4}{ms} على $T_0$.

\textbf{2.} $T = 99.4/50 = \qty{1.988}{s}$.

\textbf{3.} $L = gT^2/4\pi^2 = 9.81 \times 1.988^2/4\pi^2 \approx
\qty{0.982}{m}$.

\textbf{4.} $L = \qty{0.994}{m}$ (أي نوّاس الثواني): فأخفض
الكرة، بإطالة قدرها $\approx \qty{1.2}{cm}$.

\textbf{5.} قصير أكثر مما ينبغي، ومن ثم سريع أكثر مما ينبغي: فهي سريعة، بمقدار $(2.000 -
1.988)/1.988 \times 86400 \approx \qty{5.2e2}{s} \approx 8.7$
دقيقة في اليوم.

\textbf{6.} القمر: $\times\sqrt{9.81/1.62} = 2.46$ --- أي غير صالحة للاستعمال.
والمرصد: $\Delta T_0/T_0 = \tfrac12 \times 0.01/9.81 =
\num{5.1e-4}$: فتخسر $\approx \qty{44}{s}$ في اليوم.

\textbf{7.} $\Delta L = L\lambda\,\Delta\theta = 0.994 \times
\num{1.9e-5} \times 5.0 \approx \qty{9.4e-5}{m} \approx
\qty{0.1}{mm}$.

\textbf{8.} $T' = 2\pi\sqrt{L(1 + \lambda\Delta\theta)/g} =
T_0\sqrt{1 + \lambda\Delta\theta} \approx T_0\,(1 + \tfrac12
\lambda\Delta\theta)$.

\textbf{9.} $\Delta T_0/T_0 = \tfrac12 \times \num{1.9e-5} \times
5.0 = \num{4.8e-5}$: أي $86400 \times \num{4.8e-5} \approx
\qty{4.1}{s}$ في اليوم، متأخرةً (فالقضيب أطول، \omterm{def:g12:oscillators-and-time:oscillator}{والدور} أطول).

\textbf{10.} $\Delta\theta = \qty{-10}{K}$: $\Delta T_0/T_0 =
\num{-9.5e-5}$، فتكسب $\approx \qty{8.2}{s}$ في اليوم.

\textbf{11.} يتمدد الفولاذ والنحاس تمددين مختلفين؛ وإذا عُلّقا في اتجاهين
متناوبين، خفضت القضبان النحاسية الكرةَ بينما ترفعها الفولاذية.
وإذا قُدّرت الأبعاد بحيث تتلاشى الإزاحتان، بقي الطول الفعّال ---
ومعه المعدل --- سالمًا من الفصول.

\textbf{12.} $280/2.000 = 140$ دورًا.

\textbf{13.} $\theta_m = 2.00^\circ = \qty{0.0349}{rad}$: $E_m =
\tfrac12 \times 1.2 \times 9.81 \times 0.994 \times 0.0349^2 \approx
\qty{7.1e-3}{J}$.

\textbf{14.} في كل \omterm{def:g12:oscillators-and-time:oscillator}{دور} تتقلص \omterm{def:g12:oscillators-and-time:oscillator}{السعة} بمقدار $2^{-1/140} \approx
0.9951$، \omterm{def:g10:energy-conservation:energy}{والطاقة} بمقدار $0.9951^2 \approx 0.990$: أي نحو $1.0\%$ من
$E_m$ في \omterm{def:g12:oscillators-and-time:oscillator}{الدور}.

\textbf{15.} الدفعة $\approx 0.010 \times \num{7.1e-3} =
\qty{7.1e-5}{J}$؛ \omterm{def:g10:energy-conservation:power}{والاستطاعة} $\num{7.1e-5}/2.0 \approx \qty{3.5e-5}{W}$
--- أي عشرات الميكروواطات.

\textbf{16.} الأسبوع $604800/2.0 = \num{3.02e5}$ \omterm{def:g12:oscillators-and-time:oscillator}{دور}، ويحتاج
$\num{3.02e5} \times \num{7.1e-5} \approx \qty{21}{J}$؛ ويوفّر الثقل
$Mgh = 2.5 \times 9.81 \times 1.0 = \qty{24.5}{J}$: أي ما يكفي،
\omterm{def:g10:energy-conservation:efficiency}{بمردود} $21/24.5 \approx 88\%$.

\textbf{17.} $T = 1/\num{32768} = \qty{30.5}{\micro s}$. و$\num{32768}
= 2^{15}$: فخمس عشرة مرحلة قسمة على اثنين تحوّل \omterm{def:g10:signals-and-waves:signal}{إشارة} البلورة
إلى دقّة \qty{1}{Hz} بالضبط.

\textbf{18.} الكوارتز: $0.3/\num{2.6e6} \approx \num{1.2e-7}$.
والنوّاس: $0.5/86400 \approx \num{5.8e-6}$، أي\ $\approx
\qty{15}{s}$ في الشهر. فيفوز الكوارتز بعامل $\approx 50$.

\textbf{19.} $1/\num{3.2e15} \approx \num{3e-16}$: أي تسع رتب من
القدر دون الكوارتز. وGPS يوقّت الإشارات إلى النانوثانية ---
و\qty{30}{cm} من انتقال الضوء لكلٍّ --- فلا تفي إلا الساعات الذرّية.

\textbf{20.} تحفظ ساعة الجدار، مرمَّمةً، $\pm\qty{15}{s}$ في
الشهر في مقابل $\pm\qty{0.3}{s}$ للكوارتز: فاحتفظ بها من أجل الرنّة،
وثِق بالكوارتز من أجل القطار --- أما \omterm{def:g10:orders-of-magnitude:unit}{الثانية} نفسها فتخصّ الآن
\omterm{def:g11:fundamental-interactions:ladder}{ذرة} السيزيوم.
\end{solution}
